\documentclass[11pt]{article}

\usepackage[T1]{fontenc}
\usepackage{lmodern}
\usepackage[margin=1in]{geometry}
\usepackage{amsmath}
\usepackage{amssymb}
\usepackage{booktabs}
\usepackage{microtype}
\usepackage{xcolor}
\usepackage[hidelinks]{hyperref}
\usepackage{enumitem}

\hypersetup{
  pdftitle={Audit of the Full-Covariance Gamma-Method Implementation},
  pdfauthor={GradientFlowRGToolkit},
  pdfsubject={Scientific audit of a multivariate long-run covariance estimator}
}

\definecolor{noteBlue}{HTML}{244A73}
\definecolor{warningRed}{HTML}{8B1E2D}
\definecolor{noteGray}{HTML}{F2F4F7}

\newcommand{\code}[1]{\texttt{#1}}
\newcommand{\Cov}{\operatorname{Cov}}
\newcommand{\Var}{\operatorname{Var}}
\newcommand{\GF}{\mathrm{GF}}

\title{\textbf{Audit of the Full-Covariance Gamma-Method Implementation}\\
\large What is established, what the code computes, and what remains unvalidated}
\author{Technical note for GradientFlowRGToolkit}
\date{31 August 2026}

\begin{document}
\maketitle

\begin{center}
\fcolorbox{warningRed}{noteGray}{
  \parbox{0.88\textwidth}{
    \textbf{Current scientific status: experimental.}
    This note audits the estimator now named
    \code{ExperimentalRectangularLongRunCovariance}; it is no longer the
    primary estimator.  The adopted Bartlett/Newey--West implementation and
    its validators are documented in
    \code{long-run-covariance-methods.pdf}.
    The full long-run covariance matrix is a standard quantity already
    defined in Wolff's Gamma-method paper. That target was not invented by
    the agent that wrote the legacy code. However, the legacy implementation's
    window-selection rule, mandatory minimum window, separate-channel
    estimation, and eigenvalue clipping do not together constitute a
    literature-validated implementation of Wolff's method. A scalar
    independent-data test demonstrates substantial finite-sample bias. The
    current estimator should therefore not yet be used as authoritative
    statistical evidence.
  }
}
\end{center}

\section{The essential distinction}

Two claims must not be conflated:
\begin{enumerate}[itemsep=0.35em]
  \item \textbf{The target is justified.} For a stationary vector history,
    the covariance of the sample mean depends on the full matrix of lagged
    cross-correlations. This is established statistical theory.
  \item \textbf{The particular estimator is not yet justified.} There are
    many ways to truncate, taper, regularize, or otherwise estimate that
    infinite lag sum. The current code makes several choices that do not
    match Wolff's published practical prescription and have not otherwise
    been validated.
\end{enumerate}

The legacy implementation in \code{ContinuousBetaFunction} was generated by
a coding agent. Its presence in that codebase is provenance, not scientific
validation. The present audit therefore starts again from the primary
literature and from falsifiable synthetic tests.

\section{The full-matrix target is standard}

Let \(X_n\in\mathbb R^M\) be a stationary vector-valued Monte Carlo history,
with mean \(\mu\), and define
\[
  \Gamma_{ab}(k)
  =
  \mathbb E\left[(X_{n,a}-\mu_a)(X_{n+k,b}-\mu_b)\right].
\]
The long-run covariance, also called the zero-frequency spectral covariance,
is
\[
  \Sigma
  =
  \sum_{k=-\infty}^{\infty}\Gamma(k)
  =
  \Gamma(0)+
  \sum_{k=1}^{\infty}\left[\Gamma(k)+\Gamma(k)^{\mathsf T}\right].
\]
Under the usual Markov-chain central-limit conditions,
\[
  \sqrt N(\overline X-\mu)
  \Longrightarrow
  \mathcal N(0,\Sigma),
  \qquad
  \Cov(\overline X)\simeq\frac{\Sigma}{N}.
\]

Wolff defines the matrix-valued \(\Gamma_{\alpha\beta}(n)\), the full sum
\(C_{\alpha\beta}=\sum_t\Gamma_{\alpha\beta}(t)\), and the resulting
covariance of the means explicitly in Eqs.~(2)--(4) and (12)--(13) of
Ref.~\cite{Wolff2004}. Modern multivariate spectral-variance theory uses the
same target~\cite{VatsFlegalJones2018}. Extending beyond diagonal error bars
is therefore not a speculative extension of the Gamma method; it restores
information already present in its general formulation.

\section{What the current code computes}

For one channel, the input is an \(N\times M\) matrix. Configurations index
the rows and retained flow times index the columns. The code first centers
each column:
\[
  \delta X_{na}=X_{na}-\overline X_a.
\]
It estimates each diagonal normalized autocorrelation using
\[
  \widehat\rho_a(k)
  =
  \frac{\sum_{n=1}^{N-k}\delta X_{na}\delta X_{n+k,a}}
       {\sum_{n=1}^{N}\delta X_{na}^{\,2}}.
\]
Beginning with \(\widehat\tau_a=1/2\), it adds selected positive
autocorrelations while \(k<S\widehat\tau_a\), then chooses one matrix-wide
window
\[
  W=\left\lceil S\max_a\widehat\tau_a\right\rceil,
\]
forces \(W\geq4\) when the sample length permits, and caps it at \(N/4\).

For
\[
  G_{ab}(k)=\sum_{n=1}^{N-k}\delta X_{na}\delta X_{n+k,b},
\]
the code returns the pre-projection estimate
\[
  \boxed{
  \widehat{\Cov}(\overline X_a,\overline X_b)
  =
  \frac{1}{N^2}\left[
    G_{ab}(0)+\sum_{k=1}^{W}\{G_{ab}(k)+G_{ba}(k)\}
  \right].
  }
\]
This is a multivariate rectangular-window spectral-covariance estimator. It
does estimate off-diagonal lag corrections from data rather than inferring
them from diagonal autocorrelation times.

\section{Why this is not Wolff's published practical algorithm}

Wolff's sample autocovariance uses the lag-dependent denominator \(N-k\)
for one replica, not the constant-\(N\) convention implicit above. More
importantly, Wolff's automatic procedure defines an effective exponential
time and a function \(g(W)\), then selects the first sign change of \(g(W)\);
see Eqs.~(50)--(52) of Ref.~\cite{Wolff2004}. The code's
\(W=\lceil S\max_a\widehat\tau_a\rceil\) rule is not that procedure.

Wolff also derives a leading correction for the bias introduced by estimating
the mean from the same finite sample; see Eqs.~(32) and (49). The current code
does not apply it. A maximum over coordinate-wise autocorrelation times is
not guaranteed to control the autocorrelation time of every linear
combination. Consequently, describing this implementation simply as
``following Wolff's method'' overstates what the source supports.

\section{A decisive independent-data check}

Suppose \(X_n\) are independent scalar observations with variance
\(\sigma^2\). There is no physical autocorrelation to sum. Nevertheless,
the current code normally forces \(W=4\). Because every observation is
centered using the sample mean,
\[
  \mathbb E[G(0)]=(N-1)\sigma^2,
  \qquad
  \mathbb E[G(k)]=-\frac{N-k}{N}\sigma^2\quad(k>0).
\]
For a fixed window \(W\), the expected pre-projection estimate relative to
the true variance of the mean, \(\sigma^2/N\), is therefore
\[
  \frac{\mathbb E[\widehat{\Var}(\overline X)]}{\sigma^2/N}
  =
  1-\frac{2W+1}{N}+\frac{W(W+1)}{N^2}.
\]
With \(W=4\), the expected ratio is only \(0.60\) at \(N=20\). This failure
is structural, not a floating-point artifact.

A 20,000-replication standard-normal simulation of the actual implementation
gave:
\begin{center}
\begin{tabular}{@{}rrrr@{}}
\toprule
\(N\) & mean estimate & relative bias & negative before clipping \\
\midrule
20  & 0.621 & \(-37.9\%\) & \(14.68\%\) \\
50  & 0.823 & \(-17.7\%\) & \( 2.40\%\) \\
100 & 0.915 & \( -8.5\%\) & \( 0.16\%\) \\
400 & 0.975 & \( -2.5\%\) & approximately zero \\
\bottomrule
\end{tabular}
\end{center}
The target in every row is one after scaling by \(N\). Eigenvalue clipping
raises some negative scalar estimates to zero, explaining the small
difference from the analytic pre-projection expectation, but does not repair
the bias.

\section{Positive-semidefinite projection}

A rectangular finite lag sum need not be positive semidefinite. The current
implementation diagonalizes the estimate and replaces negative eigenvalues
by zero:
\[
  C=Q\Lambda Q^{\mathsf T},
  \qquad
  C_{\mathrm{PSD}}=Q\max(\Lambda,0)Q^{\mathsf T}.
\]
For a symmetric matrix, this is the nearest PSD matrix in Frobenius norm
\cite{Higham1988}. That is an algebraic result, not a statistical validation:
it establishes neither bias nor confidence-interval coverage, and it can
erase scientifically important low-variance directions.

The adjustment magnitude and affected directions should be recorded whenever
projection is used. A small Frobenius change is useful diagnostic evidence,
but it does not establish that the projected estimator has valid statistical
coverage.

Positive-semidefinite lag-window estimators exist. Bartlett/Newey--West
weights, for example, are PSD by construction~\cite{NeweyWest1987}. This
does not remove the need to justify bandwidth selection or validate finite
sample performance.

\section{Two additional covariance problems}

\subsection{Correlated channels are currently separated}

The raw combined \(s\) history correctly includes its covariance with the
raw \(p\) measurements used to form it. But the code invokes the estimator
separately for final \(p\), combined \(s\), and \(c\) blocks. The resulting
software objects therefore encode zero covariance between those blocks.
When the channels are measured on matched configurations, the statistically
correct unit is their jointly concatenated per-configuration vector. This
is especially important for \(p\) and combined \(s\), because the latter is
constructed partly from the former.

\subsection{High dimension implies singularity}

Any estimator expressible as \(X^{\mathsf T}KX\) from a centered
\(N\times M\) history has rank at most \(N-1\). Whenever \(M>N\), a valid PSD
covariance is therefore necessarily singular. Neither a PSD window nor
eigenvalue clipping can create information in unsupported directions.
Downstream inverse-covariance fits require an explicit dimension-reduction,
low-rank, flow-time-selection, or independently justified regularization
policy.

\section{What remains trustworthy in the processing chain}

Conditional on a valid joint covariance estimate, deterministic propagation
through
\[
  g_{\GF}^{\,2}(t)=\mathcal N(t)E(t),
  \qquad
  \beta_{\GF}(t)=-t\frac{d g_{\GF}^{\,2}}{dt}
\]
is conceptually appropriate: coupling and beta-function values should retain
their inherited correlations. The present concern is upstream---whether the
covariance supplied to that propagation is a defensible estimate.

\section{Required validation before adoption}

The estimator decision should be made explicitly rather than inherited from
the legacy agent-generated implementation. At minimum, candidates should
pass the following tests:
\begin{enumerate}[itemsep=0.35em]
  \item scalar and vector IID data recover \(\Sigma/N\) without a forced-lag
    bias;
  \item a vector AR or VAR process recovers its analytically known long-run
    covariance;
  \item linear equivariance holds:
    \(\widehat\Sigma(AX)=A\widehat\Sigma(X)A^{\mathsf T}\);
  \item selected derived scalar histories agree with a faithful projected
    Wolff calculation;
  \item cross-channel covariance is retained from matched configurations;
  \item coverage and bandwidth sensitivity are measured, not inferred from
    PSD alone; and
  \item singular high-dimensional output is represented honestly and never
    silently inverted.
\end{enumerate}

A defensible comparison would include (i) a faithful Wolff calculation for
selected scalar derived quantities and (ii) a sourced multivariate
spectral estimator, such as Bartlett/Newey--West with an explicit bandwidth
policy. Until that comparison is complete, the current rectangular estimator
plus projection should be named as an experimental comparator, not frozen in
as ``the Gamma method.''

\section{Conclusion}

The user's worry identifies a real scientific governance problem: agent-written
analysis code must not acquire authority merely by surviving in a legacy
repository. In this case, primary sources validate the desired full-matrix
quantity, which is reassuring. They do not validate the current estimator's
specific combination of heuristics. The independent-data counterexample is
sufficient to reject the present implementation as a production default.

\begin{thebibliography}{9}

\bibitem{Wolff2004}
U.~Wolff,
``Monte Carlo errors with less errors,''
\emph{Computer Physics Communications} 156 (2004) 143--153.
\href{https://arxiv.org/abs/hep-lat/0306017}{arXiv:hep-lat/0306017}.

\bibitem{VatsFlegalJones2018}
D.~Vats, J.~M. Flegal, and G.~L. Jones,
``Strong consistency of multivariate spectral variance estimators in Markov
chain Monte Carlo,''
\emph{Bernoulli} 24 (2018) 1860--1909.
\href{https://arxiv.org/abs/1507.08266}{arXiv:1507.08266}.

\bibitem{NeweyWest1987}
W.~K. Newey and K.~D. West,
``A simple, positive semi-definite, heteroskedasticity and autocorrelation
consistent covariance matrix,''
\emph{Econometrica} 55 (1987) 703--708.
\href{https://doi.org/10.2307/1913610}{doi:10.2307/1913610}.

\bibitem{Higham1988}
N.~J. Higham,
``Computing a nearest symmetric positive semidefinite matrix,''
\emph{Linear Algebra and its Applications} 103 (1988) 103--118.
\href{https://doi.org/10.1016/0024-3795(88)90223-6}{doi:10.1016/0024-3795(88)90223-6}.

\end{thebibliography}

\end{document}
